\chapter{Entrenamiento de las políticas PPO para el robot físico}
\label{cap:ppo-fisico}

Este capítulo describe cómo se formuló, entrenó, seleccionó y ejecutó la capa
correctiva PPO en el robot real entre el 16 y el 21 de octubre de 2026. Completa
el contrato de corrección residual del capítulo anterior y las campañas
preparatorias en Gazebo y MuJoCo del capítulo de resultados: aquellas usaban un
modelo nominal del actuador, mientras que las de este capítulo se entrenan con
un modelo del actuador identificado con datos del propio robot. La base teórica
de PPO está en la sección de aprendizaje por refuerzo del marco teórico.

\section{Evolución de los entornos de entrenamiento}

La capa correctiva pasó por cuatro entornos, cada uno más cercano al robot:
\begin{enumerate}
    \item \textbf{Banco cinemático en NumPy} (septiembre): observaciones,
    saturaciones y terminaciones sin dinámica ni contacto. Sirvió para validar
    la integración con cinco semillas (11, 23, 37, 53 y 71).
    \item \textbf{Gazebo} (septiembre): la política se conectó al modelo
    dinámico con contacto; ninguna escala de corrección superó a la marcha
    nominal en el criterio conjunto.
    \item \textbf{MuJoCo con modelo nominal} (septiembre y octubre):
    reentrenamiento extendido, con mejora de avance en simulación pero sin
    modelo identificado del actuador.
    \item \textbf{MuJoCo con el actuador real} (desde el 16 de octubre): entorno
    \texttt{NovaRealEnv} (programa \path{entrenar_ppo_real.py}).
    La marcha nominal es la misma que ejecuta el robot, el actuador se
    identifica con sus sensores y la observación se limita a lo que el robot
    puede medir. Es el entorno de todas las campañas de este capítulo.
\end{enumerate}

\section{Entorno de entrenamiento con el actuador real}

\subsection{Modelo físico y temporización}

El simulador es MuJoCo con paso de integración de 2~ms. La política actúa cada
30~ms (15 pasos de integración por paso de control), el mismo periodo que el
bucle de control del robot (33~ms medidos). Desde el 20 de octubre el modelo usa
la geometría medida en el robot (archivo \path{nova_sm3_real.xml}),
Tabla~\ref{tab:geometria-real}. Masas e inercias siguen siendo las del modelo de
referencia NovaSM3.

\begin{table}[H]
\centering
\caption{Geometría medida del robot e incorporada al modelo de entrenamiento.}
\label{tab:geometria-real}
\small
\begin{tabular}{lrr}
\toprule
Magnitud & Robot real (mm) & Modelo anterior (mm) \\
\midrule
Fémur (eje de cadera a eje de rodilla) & 108 & 107,95 \\
Tibia hasta la punta del pie & 155 & 153,89 \\
Desplazamiento lateral de la coxa & 35 & 38,1 \\
Separación de caderas delante--detrás & 190 & 180 \\
Separación de caderas izquierda--derecha & 77 & 120 \\
Largo del cuerpo & 265 & 230 \\
\bottomrule
\end{tabular}
\end{table}

\subsection{Marchas nominales}

La marcha nominal $\mathbf{q}^{\text{nom}}(\phi)$ de la ecuación~\eqref{eq:residual}
es una de las siguientes, con las coxas fijas en su posición de equilibrio:
\begin{itemize}
    \item \textbf{Gateo}: una pata en el aire cada vez, en el orden delantera
    izquierda, trasera derecha, delantera derecha, trasera izquierda;
    0,5~s por fase con una pausa entre fases (ciclo de 2,64~s). Fémur
    $\pm15^\circ$ y elevación de tibia de $10^\circ$.
    \item \textbf{Paso}: igual secuencia pero continua (ciclo de 2,0~s). Se
    abandonó el 20 de octubre: con las coxas fijas resultaba un gateo continuo,
    visualmente indistinguible del gateo.
    \item \textbf{Trote}: patas diagonales a la vez (delantera izquierda con
    trasera derecha y viceversa), dos en el aire y dos apoyadas; ciclo de
    1,0~s, fémur $\pm12^\circ$ y tibia $12^\circ$.
\end{itemize}
Desde el 20 de octubre ambas marchas usan una \emph{postura abierta}: el punto
neutro de los fémures delanteros se adelanta y el de los traseros se atrasa
($8^\circ$ en gateo y $5^\circ$ en trote). Con la geometría real, las patas de
un mismo lado se acercaban hasta 69~mm bajo el centro del cuerpo y llegaban a
chocar; con la apertura el hueco mínimo entre pies sube a 133~mm.

\subsection{Acción}

La acción es un vector de 12 componentes, una por articulación. Las cuatro
coxas están enmascaradas (corrección cero), por lo que la política corrige
ocho articulaciones: los cuatro fémures y las cuatro tibias. La salida de la
red se recorta a $[-1,1]$ y se escala a $\pm0{,}08$~rad ($\pm4{,}6^\circ$); la
corrección aplicada no puede cambiar más de $0{,}02$~rad ($1{,}15^\circ$) por
paso respecto a la anterior. Además, como el firmware del robot, la orden total
no puede variar más de $3^\circ$ por paso y se recorta a los topes medidos de
cada servo.

\subsection{Observación}

La observación incluye solo magnitudes medibles en el robot, porque este no
dispone de IMU. Desde la versión 6 tiene 22 componentes:
\begin{itemize}
    \item desviación de las ocho articulaciones de fémur y tibia respecto a su
    postura de referencia (8), con ruido gaussiano de $1^\circ$ y un sesgo
    aleatorio de hasta $\pm1^\circ$ por episodio, como los potenciómetros;
    \item fase de la marcha codificada como $\sin(2\pi\phi)$ y $\cos(2\pi\phi)$ (2);
    \item corrección del paso anterior normalizada por $0{,}08$~rad (12).
\end{itemize}
Las versiones 2 a 5 observaban solo tres articulaciones (fémures izquierdos y
tibia delantera izquierda), las únicas con sensor fiable en esa fecha, con 17
componentes.

\subsection{Recompensa}

La recompensa por paso es
\begin{multline}
    r_t = 1 + w_{\text{av}}\,\Delta x - w_{\text{act}}\,(\varphi^2+\vartheta^2)
    - 25\,(h-0{,}235)^2 - 0{,}25\,\lVert\mathbf{e}_q\rVert^2 - 0{,}04\,\lVert\Delta\mathbf{q}\rVert^2 \\
    - w_{\text{lat}}\,y^2 - w_{\text{giro}}\,\psi^2
    - w_{\text{suav}}\sum_j\Big(\frac{\Delta\mathbf{q}_{t,j}-\Delta\mathbf{q}_{t-1,j}}{0{,}02}\Big)^2
    - w_{\text{cox}}\,\overline{e_{\text{coxa}}^2},
    \label{eq:recompensa-real}
\end{multline}
donde $\Delta x$ es el avance del cuerpo en el paso (m); $\varphi$ y $\vartheta$,
el alabeo y el cabeceo (rad); $h$, la altura del cuerpo (m); $\mathbf{e}_q$, el
error articular respecto a la marcha nominal; $\Delta\mathbf{q}$, la corrección;
$y$ y $\psi$, el desvío lateral (m) y el giro (rad) acumulados desde el inicio;
y $\overline{e_{\text{coxa}}^2}$, la media del cuadrado de la desviación de las
coxas respecto a su orden (en grados). Además, al completar cada ciclo se resta
$w_{\text{gc}}\,(\Delta\psi_{\text{ciclo}})^2$, con el giro del ciclo en grados,
porque sin IMU la política no observa el rumbo y tendía a curvarse en caminatas
largas. Los pesos de cada versión están en la Tabla~\ref{tab:campanas-ppo}.

El término de suavidad (versión 11 en adelante) se añadió porque en el robot la
primera PPO del paso daba correcciones más rápidas de lo que los servos podían
seguir; el término de coxas, porque esa misma política sacudía el cuerpo y las
coxas, fijas en la simulación pero flexibles en el robot, se desviaban el doble
que con la marcha nominal.

\subsection{Terminación y duración de los episodios}

Un episodio termina si la altura del cuerpo sale de $[0{,}12;\,0{,}40]$~m, si
el alabeo o el cabeceo superan $0{,}6$~rad o si la observación deja de ser
finita. Si no, se trunca al completar 6 ciclos de marcha (12 en el trote, por
su ciclo más corto). Cada episodio comienza con una rampa desde la postura de
reposo hasta la postura de la marcha, durante la cual la política no actúa.

\subsection{Modelo del actuador y aleatorización}

Entre la orden y la simulación se interpone el modelo del actuador
identificado en el robot: un retardo puro (cola de 3 pasos, $\approx$99~ms), una
ganancia por articulación, una zona muerta ($\approx0{,}4^\circ$) y el límite de
$3^\circ$ por paso. Se identificó por primera vez el 16 de octubre con las
caminatas nominales registradas por los sensores del robot (retardo de 65 a
100~ms, ganancia de 0,88 a 1,0, histéresis de 0,6 a $1{,}5^\circ$) y se
actualizó después con escalones, senoides y barridos de rango.

En cada episodio se aleatorizan el retardo (2 a 4 pasos; desde el 21 de
octubre, el intervalo medido en el robot), la ganancia de cada servo
($\pm5$~\%) y el sesgo de los sensores ($\pm1^\circ$), según la estrategia de
aleatorización del dominio descrita en el marco teórico.

\section{Red neuronal e hiperparámetros}

Se usó la implementación de PPO de Stable-Baselines3~2.3.2~\cite{raffin2021sb3}
sobre PyTorch~2.3.1, con la política \texttt{MlpPolicy}: actor y crítico son
redes independientes, cada una con dos capas ocultas de 64 neuronas y activación
$\tanh$. El actor da la media de la gaussiana; las desviaciones son parámetros
aprendidos independientes del estado. La Tabla~\ref{tab:hiperparametros-ppo}
resume los hiperparámetros de las campañas desde la versión 10; en las versiones
2 a 6 no se registró la configuración de PPO en los metadatos.

\begin{table}[H]
\centering
\caption{Hiperparámetros de PPO en las campañas con el actuador real (versión 10 en adelante).}
\label{tab:hiperparametros-ppo}
\small
\begin{tabular}{lll}
\toprule
Hiperparámetro & Valor & Significado \\
\midrule
Tasa de aprendizaje & $1\cdot10^{-4}$ & paso del optimizador Adam \\
Pasos por recolección $N$ & 2048 & datos por actualización \\
Tamaño de minilote & 256 & muestras por paso de gradiente \\
Épocas por actualización $K$ & 10 & pasadas sobre cada recolección \\
Factor de descuento $\gamma$ & 0,99 & horizonte de $\approx$100 pasos (3~s) \\
Parámetro GAE $\lambda$ & 0,95 & compromiso sesgo--varianza \\
Recorte $\epsilon$ & 0,2 & cambio máximo de la razón $\rho_t$ \\
Coeficiente del crítico $c_1$ & 0,5 & peso del error de valor \\
Coeficiente de entropía $c_2$ & 0 & sin bonificación de exploración \\
Norma máxima del gradiente & 0,5 & recorte del gradiente \\
KL objetivo & 0,03 & corta las épocas si la política cambia mucho \\
Copias intermedias & cada $10^5$ pasos & para la selección posterior \\
Evaluación final & semillas 801--805 & 5 episodios deterministas \\
\bottomrule
\end{tabular}
\end{table}

\section{Campañas de entrenamiento y semillas}

La Tabla~\ref{tab:campanas-ppo} recoge todas las campañas con el actuador real.
Cada fila es una configuración entrenada con varias semillas independientes; la
columna de pasos indica el máximo alcanzado (en las continuaciones, contando los
pasos de la política de partida). Las semillas 11, 23, 37, 53 y 71 son las
congeladas en el protocolo; las demás se añadieron para ampliar la muestra. Las
continuaciones parten de la mejor política de una campaña anterior y siguen
entrenando con la configuración nueva.

\begin{table}[H]
\centering
\caption{Campañas PPO con el actuador real, 16 a 21 de octubre de 2026.}
\label{tab:campanas-ppo}
\scriptsize
\setlength{\tabcolsep}{3pt}
\begin{tabular}{p{0.1\textwidth}p{0.06\textwidth}p{0.07\textwidth}p{0.19\textwidth}p{0.08\textwidth}p{0.4\textwidth}}
\toprule
Versión & Fecha & Marcha & Semillas & Pasos & Cambio principal y resultado \\
\midrule
v2, v3 & 16/10 & gateo & 11, 23, 37, 53, 71 & 0,15 M & Primer entorno con el actuador real; 3 articulaciones observadas. Todas superan el avance nominal en simulación. \\
v4 & 16/10 & gateo & 11, 23, 37, 53, 71 & 0,2 M & Recompensa de estabilidad ($w_{\text{av}}=10$, $w_{\text{act}}=30$, $w_{\text{lat}}=20$, $w_{\text{giro}}=10$). Semilla 37: $-54$~\% de deriva y $-31$~\% de inclinación en simulación; primera ejecutada en el robot. \\
v5 & 16/10 & gateo & 11, 23, 37, 53, 71 & 0,2 M & Actuador v4 (barrido de rango). No mejora la base. \\
v6 & 19/10 & gateo & 11, 23, 37, 53, 71 & 0,2 M & Observación de las 8 articulaciones (22 componentes). \\
v10 & 19/10 & gateo, paso & 11, 23, 37, 53, 71 & 3,0--3,2 M & Aleatorización del actuador, $w_{\text{gc}}=5$, tasa $10^{-4}$. Elegidas: gateo semilla 11 (0,3 M) y paso semilla 71 (2,3 M). \\
v11 & 20/10 & paso & 11, 23, 37, 53, 61, 71, 83, 97 & 5,3 M & Suavidad ($w_{\text{suav}}=0{,}3$) y coxas ($w_{\text{cox}}=1$). Mejora a la v10, pero no a la base en el robot. \\
v12 & 20/10 & gateo & 11, 37, 41, 59 & 4,8 M & Geometría real. Elegida: semilla 41 (0,5 M). \\
trote & 20/10 & trote & 11, 23, 37, 53, 61, 71, 83, 97 & 3,0 M & Primera PPO del trote, recompensa de la v11. \\
v13 & 20/10 & gateo & 73, 89, 107, 127, 149, 163, 181, 199, 223, 251 & 3,0 M & Postura abierta $8^\circ$ y tibia $14^\circ$; pesos por defecto ($w_{\text{av}}=18$, $w_{\text{act}}=8$, sin $w_{\text{lat}}$ ni $w_{\text{giro}}$). Evaluadas solo en simulación. \\
v14 & 21/10 & gateo, trote & gateo 13, 29, 41, 43, 67; trote 19, 31, 37, 47, 89 & 3,0 M & Postura abierta; continúan copias de 0,9 M. Elegida gateo: semilla 43 (1,7 M). Ninguna de trote supera la base. \\
v15 & 21/10 & trote & 19, 211, 223, 241, 263, 281, 307 & 3--5 M & $w_{\text{av}}=20$, $w_{\text{suav}}=0{,}1$. Elegida: semilla 307 (0,1 M). \\
ciclo 1 & 21/10 & gateo, trote & 401$^{\dagger}$ & +1 M & Continúan las ganadoras en el robot. Elegida gateo: copia de 2,7 M. \\
ciclo 2 & 21/10 & gateo, trote & 401$^{\dagger}$ & +1 M & Actuador medido en el robot en la vuelta anterior. Elegida gateo: copia de 2,4 M. \\
\bottomrule
\multicolumn{6}{p{0.95\textwidth}}{\footnotesize $^{\dagger}$ Se lanzaron con las semillas 401, 409, 419 y 421, pero las cuatro ejecuciones resultaron idénticas: al continuar desde una política guardada, la biblioteca restaura la semilla almacenada en ella e ignora la indicada. Cada una de estas campañas equivale a una sola ejecución. El error se corrigió el 22 de octubre de 2026 (se impone la semilla pedida tras cargar la política).}
\end{tabular}
\end{table}

Se observaron tres patrones. Primero, la variabilidad entre semillas fue grande:
en la campaña v15 del trote, la recompensa media final de las siete semillas
osciló entre $-60$ y $+64$ con la misma configuración, lo que confirma la
necesidad de entrenar varias semillas. Segundo, la mejor copia de una campaña no
siempre es la última: varias de las elegidas tienen entre $0{,}1$ y $0{,}5$~M de
pasos, porque al seguir entrenando la política ganaba avance a costa de
inclinación o de suavidad. Tercero, las continuaciones heredan lo aprendido: las
semillas que partían de una política previa alcanzaban en pocos cientos de miles
de pasos una recompensa que las semillas nuevas tardaban millones de pasos en
lograr. La Figura~\ref{fig:curvas-semillas} muestra estos patrones en tres
campañas.

\begin{figure}[H]
\centering
\includegraphics[width=\textwidth]{Resultados/ppo_curvas_semillas.png}
\caption{Recompensa media por episodio durante el entrenamiento, una línea por
semilla, en tres campañas: gateo v14 (continuación desde copias de 0,9~M pasos),
trote v15 (la semilla 19 continúa una política previa hasta 5~M pasos) y primer
trote (8 semillas desde cero). Cada panel tiene su propia escala vertical, porque
los pesos de la recompensa difieren entre campañas. Curvas suavizadas con una
media móvil de 15 puntos.}
\label{fig:curvas-semillas}
\end{figure}

\section{Selección de copias}

Cada campaña deja una copia de la política cada $10^5$ pasos. La recompensa de
entrenamiento no basta para elegir entre ellas, porque los episodios son cortos
(6 ciclos) y sus valores dependen de los pesos de cada versión. Por eso todas
las copias se evalúan con el script \path{seleccionar_copias.py} en caminatas
simuladas largas de 60 ciclos ($\approx$3~min), con dos semillas de evaluación
(901 y 902) y en las mismas condiciones que la marcha nominal. Una copia es
\emph{válida} si avanza más que la marcha nominal, no se desvía más del límite
fijado y no se cae. Entre las válidas se elige la de menor puntuación
\begin{equation}
    P = \frac{d_y}{0{,}1\,\text{m}} + \varphi_{\text{rms}} + \vartheta_{\text{rms}}
    - \frac{x}{x_{\text{nom}}} + 2\left(\frac{c_{\text{rms}}}{c_{\text{rms,nom}}}-1\right) + S,
    \label{eq:puntuacion-seleccion}
\end{equation}
con $d_y$ el desvío lateral final, $\varphi_{\text{rms}}$ y $\vartheta_{\text{rms}}$
el alabeo y el cabeceo RMS en grados, $x$ el avance, $c_{\text{rms}}$ la
desviación RMS de las coxas y $S$ el cambio cuadrático medio de la corrección
(suavidad); el subíndice nom indica la marcha nominal. La copia ganadora se
guarda en \path{ppo_v10_robot/} junto con sus metadatos (marcha, amplitudes,
apertura, versión del actuador y resultado de la evaluación), que el programa
del robot lee para reproducir exactamente la misma configuración.

\section{Ejecución en el robot}

\subsection{Programa de control y realimentación}

En el robot, la política se ejecuta con el programa de control
\path{ejecutar_politica_robot.py} (en el PC durante los ensayos de octubre). En cada paso de 30~ms el programa calcula la
referencia nominal, construye la observación con los sensores, evalúa la red
(media de la gaussiana, sin muestreo), aplica las saturaciones, convierte los
ángulos del modelo a grados de servo con la calibración de cada articulación y
envía las doce órdenes a la interfaz de E/S (modo de transmisión continua del
firmware), que responde con todas las lecturas de sensores.

Cada articulación usa una cadena de fuentes de realimentación: su sensor
principal (potenciómetro leído a través del multiplexor analógico, o encoder
magnético AS5600); si falla, su sensor de respaldo; y si tampoco hay, la orden
retrasada 99~ms como estimación, Tabla~\ref{tab:sensores-articulaciones}. Las
cámaras, que identifican cada pieza de la pata por el color de su cinta (fémur
rojo, tibia verde, pie amarillo), se usan como tercera fuente, aunque solo en el
análisis posterior.

\begin{table}[H]
\centering
\caption{Fuentes de realimentación por articulación al 21 de octubre de 2026.}
\label{tab:sensores-articulaciones}
\footnotesize
\begin{tabular}{llll}
\toprule
Articulación & Principal & Respaldo & Estado \\
\midrule
Fémur 1 (delantera derecha) & potenciómetro C8 & encoder ENC\_F1 & operativo \\
Fémur 2 (delantera izquierda) & potenciómetro C5 & --- & operativo \\
Fémur 3 (trasera derecha) & potenciómetro C9 & --- & operativo \\
Fémur 4 (trasera izquierda) & potenciómetro C6 & encoder ENC\_F4 & operativo \\
Tibia 1 & encoder ENC\_T1 & --- & sin lecturas: orden retrasada \\
Tibia 2 & encoder ENC\_T2 & --- & operativo \\
Tibia 3 & potenciómetro C11 & --- & sin señal: orden retrasada \\
Tibia 4 & potenciómetro C10 & --- & operativo \\
\bottomrule
\end{tabular}
\end{table}

\subsection{Correcciones necesarias para que el robot reproduzca la simulación}

Comparar el robot con la simulación mediante las cámaras reveló diferencias que
ningún entrenamiento podía compensar y que se corrigieron en la interfaz
robot--modelo:
\begin{itemize}
    \item \textbf{Signo de las tibias izquierdas.} Las tibias 2 y 4 están
    montadas en espejo: al pedirles levantar el pie lo bajaban. Se invirtió su
    signo; afectaba a todas las caminatas anteriores al 21 de octubre.
    \item \textbf{Postura base.} En la postura de reposo, la simulación tenía las
    rodillas dobladas unos $48^\circ$ y el robot casi rectas. Un ajuste
    automático en bucle, que mide con las cámaras el eje de la cinta de cada
    fémur y cada tibia y corrige el 70~\% de la diferencia por vuelta, dejó los
    fémures a menos de $2^\circ$ de la simulación y las tibias a unos $10^\circ$
    (antes, hasta $50^\circ$).
    \item \textbf{Topes del firmware.} Las correcciones de postura llevaban a
    las tibias 2 y 4 más allá de sus topes de seguridad, donde quedaban fijas.
    El programa respeta ahora esos topes y reserva margen para levantar el pie.
    \item \textbf{Escala de los sensores.} En la postura nueva los
    potenciómetros daban $3{,}2$ a $3{,}7$ cuentas por grado frente a las
    $2{,}5$ de la calibración anterior; la cámara confirmó que el motor sí se
    movía lo ordenado (ganancia $\approx1{,}08$), de modo que se recalibró la
    escala de los sensores y se mantuvo la ganancia del actuador en 1.
    \item \textbf{Temporización.} Con el PC cargado por el análisis de video y
    los entrenamientos, hasta 714 pasos de una caminata de 3~min llegaban con más
    de 40~ms de retraso. Se reservaron dos núcleos para el bucle de control, la
    red se evalúa con un único hilo y los análisis corren con prioridad mínima;
    con el PC entrenando a la vez, el número de pasos tardíos bajó a cero.
\end{itemize}

\section{Ciclo de mejora en la plataforma física}

Para que el robot no solo evaluara las políticas sino que también mejorara la
simulación, se implementó un ciclo automático (\path{ciclo_mejora.py}) con
cuatro pasos por vuelta:
\begin{enumerate}
    \item \textbf{Prueba en el robot}: la marcha nominal y las políticas
    candidatas de gateo y trote, dos repeticiones de 3~min cada una, en orden
    alternado entre rondas y con pausas de 2~min para enfriar los servos.
    \item \textbf{Puntuación con los sensores}: para cada condición se promedia
    el error RMS de seguimiento de fémures y tibias y la desviación media de las
    coxas, en grados. Gana la de menor puntuación.
    \item \textbf{Actualización del modelo}: con las caminatas de la vuelta se
    vuelven a medir, por articulación, el retardo, la histéresis y el error que
    el modelo no explica (\path{actualizar_modelo_real.py}). Se genera un modelo
    del actuador nuevo, con el intervalo de retardo medido para la
    aleatorización.
    \item \textbf{Reentrenamiento}: la política ganadora de cada marcha se sigue
    entrenando durante $10^6$ pasos con ese modelo, y su mejor copia se
    convierte en candidata de la vuelta siguiente. Se previeron cuatro semillas
    por marcha (401, 409, 419 y 421), pero por el error descrito en la
    Tabla~\ref{tab:campanas-ppo} las cuatro dieron la misma ejecución, de modo
    que en las vueltas realizadas cada marcha se reentrenó, en la práctica, con
    una sola semilla.
\end{enumerate}
De este modo la selección se decide con el robot real y no con la simulación, y
las discrepancias medidas en el robot se trasladan al entrenamiento. Una
limitación de este método es que la puntuación no mide la locomoción
propiamente dicha: el robot está suspendido, por lo que se evalúan el
seguimiento de las órdenes y la estabilidad de las coxas, no el avance ni la
inclinación del cuerpo sobre el suelo.
